\chapter*{ЗАКЛЮЧЕНИЕ}
\addcontentsline{toc}{chapter}{ЗАКЛЮЧЕНИЕ}

В настоящей работе проведено систематическое исследование когнитивных предубеждений в девяти современных больших языковых моделях на трёх классических психологических парадигмах, охватывающих вероятностное суждение (задача Линды), дедуктивное рассуждение (задача выбора карточек Уэйсона) и индуктивное открытие правил (задача 2-4-6). Разработан воспроизводимый экспериментальный фреймворк с оригинальными датасетами, контролем контаминации и системой непрерывных метрик.

Основные результаты работы:

\begin{enumerate}
  \item \textbf{Когнитивные предубеждения в LLM являются широко распространённым, но не абсолютно универсальным феноменом.} Токен-bias в форме ошибки конъюнкции выявлен у восьми из девяти моделей (delta\_bias от $+0{,}028$ до $+0{,}565$, McNemar $p < 10^{-5}$). Предубеждение подтверждения в дедуктивном рассуждении воспроизведено у всех девяти моделей (совокупный индекс bias от 0,4\% до 53,6\%). В индуктивной задаче CTR $\geq 0{,}629$ у всех non-reasoning моделей. Единственное исключение~--- \texttt{qwen3.5-397b-a17b}, устойчивая к конъюнктивной ошибке, но демонстрирующая выраженное предубеждение подтверждения в задачах Уэйсона и 2-4-6.

  \item \textbf{Метрика Severity of Bias обнаруживает предубеждение за пределами бинарной классификации.} На примере \texttt{claude-sonnet-4.6} (единственная бинарная ошибка из 520 пар, но delta\_bias $= +0{,}105$, Cohen's $d = 1{,}619$) показано, что токен-предубеждение существует на спектре и не сводится к дихотомии «ошибается / не ошибается». Модели совершают конъюнктивные ошибки с уверенностью 72--99 из 100, что опровергает гипотезу о неуверенных догадках.

  \item \textbf{Масштаб модели не является монотонным предиктором устойчивости к предубеждениям.} Внутрисемейный парадокс зафиксирован в обоих протестированных семействах: в Gemma и Qwen меньшая модель демонстрирует более низкий уровень bias в вероятностной задаче, при этом направление эффекта инвертируется в логической задаче. Это указывает на мультидимензиональную природу когнитивных предубеждений: уязвимость к разным типам предубеждений не определяется единым латентным фактором.

  \item \textbf{Интенсивность предубеждений возрастает по спектру когнитивной сложности.} От вероятностного суждения (ошибка устраняется семантически нейтральным хобби) к дедукции (EMA лучших моделей не достигает 1{,}0) и далее к индукции (максимальный SR $= 0{,}28$ у non-reasoning моделей). Индуктивное открытие правил остаётся наиболее уязвимой областью.

  \item \textbf{Reasoning обеспечивает качественный скачок в индуктивной задаче.} SR \texttt{gpt-5.2} reasoning + adaptive ($0{,}48$) вдвое превышает лучший non-reasoning результат ($0{,}28$) и является единственной комбинацией, приближающейся к фальсификационной стратегии (CTR $= 0{,}559$). При этом prompt engineering (CoT, adaptive) не даёт сопоставимого эффекта и может ухудшать результаты у моделей с недостаточной базовой компетентностью.

  \item \textbf{Деонтическая структура задачи фасилитирует логический вывод независимо от знакомости содержания.} Условие «фантастических социальных контрактов» даёт EMA, сопоставимое со знакомыми контрактами (медианный прирост от абстрактного нейтрального условия $+0{,}072$), что подтверждает гипотезу прагматических схем и является оригинальным вкладом в исследования LLM.

  \item \textbf{Демографические эффекты присутствуют, но не являются универсальными.} Расовые различия статистически значимы только у \texttt{glm-4.7-flash}. Средние fallacy rate различаются менее чем на 1{,}5~п.п. между расовыми группами. Возраст не влияет на интенсивность токен-bias.
\end{enumerate}

Разработанный экспериментальный фреймворк, включая оригинальные датасеты, метрики и статистический протокол, опубликован в открытом доступе и может применяться для профилирования когнитивных уязвимостей LLM перед их развёртыванием в системах поддержки принятия решений.

Таким образом, проведённое исследование вносит вклад в понимание природы когнитивных предубеждений в языковых моделях, предлагает методологию их количественной оценки и формулирует практические рекомендации по митигации для инженеров и исследователей.
